%% 自動生成（gen.py）
\bktocchap{2-1}{「かぞえる\R{数}{かず}」と「はかる\R{数}{かず}」}{「5」だけでは、\R{何}{なん}のことかわからない}{7}{1}{1}{18}{C}
\bktocchap{2-2}{マイナスの\R{数}{かず}と\R{数直線}{すう|ちょく|せん}}{0\,℃より\R{寒}{さむ}いときは、どう\R{書}{か}く？}{6}{1}{2}{38}{C}
\bktocchap{2-3}{＋ － × ÷ は、\R{何}{なに}を\R{表}{あらわ}しているの？}{かけ\R{算}{ざん}の\R{答}{こた}えは、\R{長方形}{ちょう|ほう|けい}の\R{広}{ひろ}さ？}{1,2}{1}{1}{60}{C}
\bktocchap{2-4}{\R{分数}{ぶん|すう}と\R{小数}{しょう|すう}の\R{計算}{けい|さん}}{ピザ2\R{枚}{まい}を3\R{人}{にん}で\R{分}{わ}けると、1\R{人分}{にん|ぶん}は\R{何枚}{なん|まい}？}{1}{1}{2}{76}{C}
\bktocchap{2-5}{\R{比}{ひ}・\R{割合}{わり|あい}・\R{百分率}{ひゃく|ぶん|りつ}}{「100\,\%」をこえることはある？}{7}{1}{2}{104}{C}
\bktocchap{2-6}{はかった\R{数}{かず}には「ずれ」がある}{\R{同}{おな}じ\R{物}{もの}を10\R{回}{かい}はかると、ぜんぶ\R{同}{おな}じになる？}{7}{1}{2}{122}{C}
